%% ISMB 2026 Submission Cover Letter
\documentclass[11pt]{letter}
\usepackage[utf8]{inputenc}
\usepackage[T1]{fontenc}
\usepackage{times}
\usepackage[margin=2.5cm]{geometry}

\signature{Mingyu Sun}
\address{Mingyu Sun \\ Affiliation to be completed \\ Email: [email]}

\begin{document}

\begin{letter}{Programme Committee \\ ISMB 2026}

\opening{Dear Programme Committee Members,}

I am pleased to submit the manuscript entitled \textbf{``Systematic benchmarking of
functional module partitioning in GO pathway subnetworks: multi-resolution community
detection versus hierarchical co-expression clustering''} for consideration at ISMB 2026.

\vspace{6pt}
\noindent\textbf{Summary.} Weighted gene co-expression network analysis (WGCNA) is the
\textit{de facto} standard for partitioning genes into functional modules, yet its
application to single-pathway subnetworks rests on four assumptions that have not been
systematically tested. We assembled a paired benchmark of 11 expression cohorts, seven GO
biological-process pathways and 31 partitioning algorithms drawn from eight methodological
families, with all algorithms operating on an identical topological overlap matrix. Each
algorithm--pathway--cohort triple was evaluated across six orthogonal layers (structure,
functional coherence, stability, robustness, biological validity and downstream clinical
prediction), complemented by three synthetic benchmarks with known ground truth and two
null models.

\vspace{6pt}
\noindent\textbf{Principal findings.} Multi-resolution community detection produced
functionally more coherent modules than hierarchical clustering (median $z = 0.768$ versus
$-0.075$; paired Wilcoxon $q = 5.3\times10^{-7}$) while simultaneously improving modularity,
and recovered nested functional sub-structures in synthetic benchmarks where hierarchical
cutting merged them (ARI $1.000$ versus $0.908$). Downstream clinical prediction concurred
(AUC $0.989$ versus $0.898$).

\vspace{6pt}
\noindent\textbf{Why this is of interest to ISMB.} Beyond the algorithmic comparison, our
study delivers three methodological results that we believe are of general relevance to the
computational biology community:

\begin{enumerate}
\item \textbf{Enrichment $p$-values are size-biased.} Within a functionally coherent
pathway, enrichment-based module scores are inflated for any partition; we propose a
size-matched, pathway-internal null that removes both the size and background biases.
\item \textbf{Training an embedding model on community labels and then using it to validate
those labels is circular.} We show how self-supervised edge prediction removes this leakage
and yields a genuinely independent test.
\item \textbf{Modularity does not exceed degree-preserving null models} on these networks
(mean $z = -5.54$), which argues against using modularity as evidence that inferred modules
reflect organisation beyond what the degree distribution implies.
\end{enumerate}

We further report two negative results without suppression (the null-model finding above and
the only partial concordance between self-supervised embeddings and Leiden communities), in
line with best practice for computational benchmarking.

\vspace{6pt}
\noindent\textbf{Declarations.} This work is original, has not been published elsewhere and
is not under consideration by another venue. All datasets are publicly available from GEO.
Analysis code, processed matrices and all evaluation tables are provided as a reproducible
pipeline. The author declares no competing interests.

\vspace{8pt}
\noindent We would be grateful for the committee's consideration and look forward to the
reviewers' comments.

\closing{Sincerely,}

\end{letter}

\end{document}
